% 第2章：形式化建模与安全性质（精简定稿）
\subsection{攻击者模型与研究假设}
\label{subsec:formal-adversary}

本文采用 Dolev--Yao 主动网络攻击者：攻击者可读取公开材料，拦截、重组、替换及重放消息，但只能使用符号模型允许的构造与消去规则。普通 KEM 封装可由公开公钥执行，因此攻击者能够构造新的 $ct_l$、$ct_k$；在没有相应 split-KEM 秘密材料时，不能凭空生成匹配的诚实 $ct_s$，但可以复用已观察到的分量。主要无泄露实验不开放秘密状态，独立泄露实验仅按各自模型解释。

批次模型还允许攻击者在抽象输入边界选择条目及其槽位，包括重复投递已观察到的有效消息。这是\emph{模型中的批次组合能力}，并不证明攻击者能够操控真实部署的批次构造组件。下述 ProVerif 与 Tamarin 模型有不同的事件和密码学抽象，彼此结果不构成一条自动衔接的攻击轨迹。

\subsection{ProVerif 核心模型}
\label{subsec:proverif-core}

ProVerif 模型保留式\eqref{eq:core-transcript}中的三密文消息、预密钥签名检查、KEM/split-KEM 解封装和密钥派生，以公开信道表示网络。发送方从网络获得声称的接收方预密钥；接收方使用本地保存的接收方预密钥与解封装秘密。模型采用\emph{单次接收抽象}，不包含 \batchreceive 的批次槽位与关闭状态。

双方导出密钥时分别记录 $\mathsf{SenderKey}(A,B,sid,k)$ 与 $\mathsf{ReceiverKey}(B,A,sid,k)$；发送方完成和接收方完成分别记录
\begin{equation}
 \mathsf{SendDone}(A,B,sid,k),\qquad
 \mathsf{RecvDone}(B,A,sid,k).
 \label{eq:core-events}
\end{equation}
发送方的完成事件位于核心消息输出之后；原始模型的接收完成事件位于成功解封装及密钥派生之后。另以分量发送和分量接受事件检查 split-KEM 来源。上述事件是\emph{模型观察点}，$\mathsf{RecvDone}$ 不表示真实软件已安装会话。独立的 HMAC 变体在接收完成前增加确认检查；具体机制及验证对照留待第3章。

\subsection{Tamarin 批次模型}
\label{subsec:tamarin-batch}

Tamarin 使用多集重写规则表达线性状态和事件先后关系\cite{meier2013tamarin}。接收状态实例以 $rst$ 标识，批次以 $bid$ 标识，同批槽位以 $idx$ 区分；输入收集、封闭、逐槽处理及批次完成或失败由线性事实约束。这些阶段是本文的\emph{符号建模分解}，并非原协议新增加的算法接口。V6/V7 用于验证状态和生命周期约束，其中 V7 为固定四槽模型；批内重放与后续加固主要采用固定两槽模型。

在历史重放模型中，发送完成事件为 $\mathsf{SenderSession}(A,B,m,sid,k)$；接收发生事件为
\begin{equation}
 \mathsf{ReceiverAccept}(B,A,bid,idx,rst,m,sid,k).
 \label{eq:tamarin-accept}
\end{equation}
模型以持久的诚实会话来源事实抽象成功解封装与参数匹配，同时保留不同槽位对同一来源的重复读取可能性；它\emph{不独立重演} ProVerif 中的全部密码检查。HMAC-only 模型增加确认消息及对应接受事件；影响模型另引入有条件的输出消费与符号安装层。独立的 RQ-v2 准入原型只比较消息和参与方坐标，其发送发生标识并非协议会话标识 $sid$。这些理论的同名事件不能直接互换。

\subsection{安全性质与验证分工}
\label{subsec:formal-properties}

本文将安全目标分层，而不视为可互相替代的单一认证性质。P0-S 检查发送方或接收方密钥事件中的 $k$ 是否可由攻击者推导；P0-O 检查已接受 split-KEM 分量是否具有匹配发送来源。P1 要求每个接收完成具有更早且参数完全一致的发送完成，形式为
\begin{equation}
 \mathsf{RecvDone}(B,A,sid,k)@r\;\Longrightarrow\;
 \exists s<r:\mathsf{SendDone}(A,B,sid,k)@s.
 \label{eq:property-p1}
\end{equation}
该性质为\emph{非单射对应关系}，即同一个发送完成可以匹配多次接收；它不能保证重放防护。

P2 针对有界批次中的接收\emph{发生}。令 $\theta=(A,B,m,sid,k)$、$\beta=(bid,rst)$，$R(\theta,\beta,idx)@r$ 表示式\eqref{eq:tamarin-accept}中的接收事件，且发送来源的存在性与唯一性分别受相应引理约束。批内发生单射性的核心限制为
\begin{equation}
 R(\theta,\beta,i)@r_1\land R(\theta,\beta,j)@r_2
 \;\Longrightarrow\;r_1=r_2.
 \label{eq:property-p2}
\end{equation}
式\eqref{eq:property-p2}仅限制同一匹配消息及同一 $(bid,rst)$ 内的重复接受，不推出不同消息所属参与方必然不同。P3 则在显式的 $C_{\mathrm{install}}\text{-v2}$ 消费假设下检查同一匹配发送是否引起重复符号安装；具体输出来源和安装唯一性定义留到第4章。参与方级准入另由式\eqref{eq:party-distinct}刻画，不能由 P2 或精确消息去重直接替代。

\begin{center}
\begin{minipage}{\linewidth}
\centering
\captionof{table}{安全性质与验证模型的对应关系}
\label{tab:security-goal-map}
\small
\begin{tabularx}{\linewidth}{@{}lYlY@{}}
\toprule
性质 & 检查目标 & 工具 & 适用边界\\
\midrule
P0-S & 会话密钥不可推导 & ProVerif & 无泄露符号基线\\
P0-O & split-KEM 分量有来源 & ProVerif & 分量级对应\\
P1 & 接收完成存在匹配发送 & ProVerif & 完整参数、非单射\\
P2 & 同批相同消息至多接受一次 & Tamarin & 固定两槽、同状态\\
P3 & 同源输出的安装唯一性 & Tamarin & 条件化消费模型\\
参与方互异 & 同批不同槽位关联不同参与方 & Tamarin & 独立准入原型\\
\bottomrule
\end{tabularx}
\end{minipage}
\end{center}

表\ref{tab:security-goal-map}概括两类工具的分工。ProVerif 负责单次核心过程的机密性及对应关系；Tamarin 负责有界状态化批次中的重复接受、条件化影响和准入比较。存在性见证与全称安全性质分别报告；后续各章仅在对应模型的范围内解释证明和反例，不将不同理论的事件拼接为未经验证的联合执行。
